\scanpage{570}
\mspage{479}
\uline{Proposition} 4.16.7.2. \unsure{Soient|Soit} $M$ une catégorie
accessible, $x$ un objet de $M$, et $ {}_xM$ la
sous-catégorie fermée de $M$ engendrée
par $x$, formée des $y\in\operatorname{Ob}M$ tels que
$\operatorname{Hom}(x,y)\ne\varnothing$, i.e. \deleted{\unsure{de}} l'image essentielle
de $ {}_{x\backslash}M\mathrel{\unsure{-|\longrightarrow}}M$. Alors $ {}_xM$ est une
sous-catégorie de $M$ accessible dans $M$.

En effet, c'est l'image inverse de la
sous-catégorie $\operatorname{Ens}^{*}$ accessible dans $\unsure{\operatorname{En}|\operatorname{Ens}}$,
par le foncteur représentable (donc
accessible)
\[
h_x:y\longmapsto\operatorname{Hom}(\unsure{xy|x,y}):M\mathrel{-}(\operatorname{Ens}),
\]
donc est accessible \inserted{dans $M$} en vertu de 4.13.\ldots

\uline{\deleted{Proposition}\inserted{Corollaire}} 4.16.7.3. Soit $M$ une catégorie
accessible, $\underline S$ une sous-catégorie fermée.
Pour que $\underline S$ soit accessible, il f. et il
s. qu'il existe une \uline{petite} partie
\inserted{$J$} de $\underline S$, telle que $\operatorname{Ob}\underline S=\bigcup_{x\in J}\deleted{\unsure{\mathcal E}}\operatorname{Ob}({}_xM)$.

En effet, on a
\[
\operatorname{Ob}\underline S=\bigcup_{x\in\operatorname{Ob}\underline S}\operatorname{Ob}({}_xM),
\]
i.e. $\underline S$ apparaît comme réunion des
% archive sheet number: 569
